Das Ausgleichsgefäss muss die ganze Volumenzunahme $\Delta V = \VAgO$ des Wassers aufnehmen. Der Ausdehnungskoeffizient von Wasser hängt stark von der Temperatur ab, deshalb rechnet man mit den Dichten $\rho_1 = \rhAO$ bei \TAzO\ und $\rho_2 = \rhBO$ bei \TBzO\ (FoTaBe S.~192). Die Masse des Wassers bleibt gleich, $\rho_1 V_0 = \rho_2 V_2$, also ist $V_2 = V_0\,\rho_1/\rho_2$ und
\begin{align}
 \Delta V &= V_2 - V_0 = V_0\left(\frac{\rho_1}{\rho_2}-1\right)
\end{align}
Aufgelöst nach dem Volumen bei \TAzO:
\begin{align}
 V_0 &= \VWbF = \frac{\VAg}{\rhA/\rhB-1} = \VWb \approx \result{\VWblP{0}{2}}
\end{align}

\emph{Bemerkung:} Mit dem Ausdehnungskoeffizienten $\gamma = \gaWaO$ aus der Tabelle (FoTaBe S.~191), der bei \TtbO\ gilt, und $\Delta V = \gamma\,V_0\,\Delta\vartheta$ mit $\Delta\vartheta = \dTzF = \dTz$ ergäbe sich $V_0 = \VWaF = \VWlP{0}{2}$. Weil sich Wasser bei höheren Temperaturen stärker ausdehnt, ist dieser Wert zu gross: Mit so viel Wasser wäre das Gefäss bereits zu klein.
